\section{Sujet n\textsuperscript{o}\,13}
\noindent\begin{minipage}{0.5\linewidth}\footnotesize Office du Baccalauréat du Cameroun\end{minipage}%
\hfill\begin{minipage}{0.45\linewidth}\footnotesize\raggedleft Examen : \textbf{Baccalauréat} · Série : \textbf{C/E}\\
Épreuve : \textbf{Mathématiques} · Durée : \textbf{4\,h} · Coef. : \textbf{7}\end{minipage}
\par\smallskip{\color{onbo}\rule{\linewidth}{1pt}}

\subsection*{Partie A — Évaluation des ressources \hfill (15 points)}

\noindent\textbf{Exercice 1.} \hfill\textit{(3 points)}\\
Une boulangerie industrielle de Yaoundé produit ses baguettes au moyen de trois fours
$A$, $B$ et $C$, qui fournissent respectivement \SI{50}{\percent}, \SI{30}{\percent} et
\SI{20}{\percent} de la production journalière. La proportion de baguettes défectueuses
est de \SI{2}{\percent} pour le four $A$, \SI{4}{\percent} pour le four $B$ et
\SI{5}{\percent} pour le four $C$. On prélève une baguette au hasard dans la production
du jour.
\begin{enumerate}[label=\arabic*)]
\item Représenter la situation par un arbre pondéré, puis calculer la probabilité que la
baguette prélevée soit défectueuse. \hfill\textit{(1,5 pt)}
\item La baguette prélevée est défectueuse. Quelle est la probabilité qu'elle provienne du four $A$\,? \hfill\textit{(1 pt)}
\item Calculer de même la probabilité qu'elle provienne du four $C$. \hfill\textit{(0,5 pt)}
\end{enumerate}

\smallskip
\noindent\textbf{Exercice 2.} \hfill\textit{(3 points)}\\
La boulangerie conditionne ses pains spéciaux dans des cartons de deux modèles : un
grand modèle de $11$ pains et un petit modèle de $7$ pains. Une commande porte sur un
total de $100$ pains exactement, sans pain en surplus.
\begin{enumerate}[label=\arabic*)]
\item À l'aide de l'algorithme d'Euclide ou de l'identité de Bézout, déterminer un couple
$(x_0\,;\,y_0)\in\Z^2$ tel que $11x_0-7y_0=1$. \hfill\textit{(1 pt)}
\item Résoudre dans $\Z^2$ l'équation $(E):\ 11x+7y=100$. \hfill\textit{(1,25 pt)}
\item En déduire le nombre de grands cartons et de petits cartons nécessaires, sachant que
$x$ et $y$ doivent être des entiers \emph{positifs}. \hfill\textit{(0,75 pt)}
\end{enumerate}

\smallskip
\noindent\textbf{Exercice 3.} \hfill\textit{(4 points)}\\
Soit $f$ la fonction définie sur $\R$ par $f(x)=x\,\mathrm{e}^{-x}$, de courbe
représentative $(\mathcal{C})$ dans un repère orthonormé d'unité \SI{2}{\centi\metre}.
\begin{enumerate}[label=\arabic*)]
\item Déterminer les limites de $f$ en $-\infty$ et $+\infty$ ; interpréter graphiquement celle en $+\infty$. \hfill\textit{(1 pt)}
\item Montrer que $f'(x)=(1-x)\mathrm{e}^{-x}$, puis dresser le tableau de variations de $f$. \hfill\textit{(1,25 pt)}
\item À l'aide d'une intégration par parties, montrer que pour tout réel $\lambda>0$,
$\displaystyle\int_0^{\lambda}x\,\mathrm{e}^{-x}\,\dd x=1-(\lambda+1)\mathrm{e}^{-\lambda}$. \hfill\textit{(1 pt)}
\item En déduire l'aire, en \si{\centi\metre\squared}, du domaine limité par $(\mathcal{C})$,
l'axe des abscisses et les droites $x=0$ et $x=1$. \hfill\textit{(0,75 pt)}
\end{enumerate}

\smallskip
\noindent\textbf{Exercice 4.} \hfill\textit{(5 points)}\\
Dans le plan, on considère un triangle $ABC$ tel que $AB=4$, $AC=4$ et l'angle en $A$ soit
droit. On munit le plan du repère orthonormé $(A\,;\,\vec\imath,\vec\jmath)$ tel que
$B(4\,;\,0)$ et $C(0\,;\,4)$.
\begin{enumerate}[label=\arabic*)]
\item Soit $G$ le barycentre du système $\{(A,2)\,;\,(B,1)\,;\,(C,1)\}$. Déterminer les
coordonnées de $G$. \hfill\textit{(1 pt)}
\item Pour tout point $M$ du plan, on pose $f(M)=2MA^2+MB^2+MC^2$. Démontrer que pour tout
$M$, $f(M)=4\,MG^2+24$. \hfill\textit{(2 pt)}
\item Déterminer et caractériser l'ensemble $(\Gamma)$ des points $M$ du plan tels que
$f(M)=40$. \hfill\textit{(1,5 pt)}
\item Le point $A$ appartient-il à $(\Gamma)$\,? Justifier. \hfill\textit{(0,5 pt)}
\end{enumerate}

\subsection*{Partie B — Évaluation des compétences \hfill (5 points)}

\noindent\textit{Madame Mballa gère une boulangerie industrielle à Yaoundé et sollicite
ton aide, en tant qu'élève de Terminale~C, pour trois décisions de gestion. Le coût
total de production journalier, lorsqu'elle fabrique $x$ centaines de baguettes
($x>0$), est de $C(x)=x^2+3600$ (en milliers de francs CFA), et elle veut produire la
quantité qui rend le \emph{coût moyen par centaine} minimal. Par ailleurs, elle place
ses bénéfices : un capital initial de \SI{500000}{} francs CFA placé au taux annuel de
\SI{6}{\percent} à intérêts composés. Enfin, sur sa chaîne d'emballage,
\SI{4}{\percent} des paquets présentent un défaut ; un client en achète $5$ au hasard.}

\begin{enumerate}[label=\textbf{Tâche \arabic*.},leftmargin=2.4em]
\item Déterminer la quantité $x$ (en centaines de baguettes) qui \emph{minimise le coût
moyen} $C_m(x)=\dfrac{C(x)}{x}$, ainsi que ce coût moyen minimal. \hfill\textit{(1,5 pt)}
\item Déterminer le nombre d'années au bout duquel le capital placé aura \emph{au moins doublé}. \hfill\textit{(1,5 pt)}
\item Déterminer la probabilité qu'\emph{au moins un} des $5$ paquets achetés soit défectueux. \hfill\textit{(1,5 pt)}
\end{enumerate}
\par\smallskip{\footnotesize\itshape Présentation générale : 0,5 point.}

% ════════════════════════════════════════════════════
\corrige{du sujet 13}

\noindent\textbf{Partie A — Exercice 1.}
1) On note $A,B,C$ les événements « la baguette vient du four $A,B,C$ » et $D$
« la baguette est défectueuse ». L'arbre porte $P(A)=0{,}5$, $P(B)=0{,}3$, $P(C)=0{,}2$
et $P_A(D)=0{,}02$, $P_B(D)=0{,}04$, $P_C(D)=0{,}05$. Par la formule des probabilités totales :
$P(D)=0{,}5\times0{,}02+0{,}3\times0{,}04+0{,}2\times0{,}05=0{,}01+0{,}012+0{,}01=0{,}032=\frac{4}{125}$.
2) $P_D(A)=\dfrac{P(A\cap D)}{P(D)}=\dfrac{0{,}01}{0{,}032}=\dfrac{10}{32}=\dfrac{5}{16}\approx0{,}3125$.
3) $P_D(C)=\dfrac{0{,}01}{0{,}032}=\dfrac{5}{16}\approx0{,}3125$ (et $P_D(B)=\frac{0{,}012}{0{,}032}=\frac38$,
la somme valant bien $\frac5{16}+\frac38+\frac5{16}=1$).

\noindent\textbf{Exercice 2.}
1) Euclide : $11=1\times7+4$, $7=1\times4+3$, $4=1\times3+1$. En remontant :
$1=4-3=4-(7-4)=2\times4-7=2(11-7)-7=2\times11-3\times7$. Donc $11\times2-7\times3=1$ :
un couple est $(x_0\,;\,y_0)=(2\,;\,3)$.
2) En multipliant par $100$ : $11\times200-7\times300=100$, soit une solution
particulière $(200\,;\,-300)$ de $(E)$. Par différence, $11(x-200)=-7(y+300)$, donc
$11(x-200)=7\,(-(y+300))$ ; comme $11\wedge7=1$, le théorème de Gauss donne $7\mid x-200$,
d'où $x=200+7k$ et alors $y=-300-11k$, $k\in\Z$. (On vérifie : $11(200+7k)+7(-300-11k)=2200-2100=100$.)
$\boxed{\,S=\{(200+7k\,;\,-300-11k),\ k\in\Z\}\,}$.
3) Conditions $x\ge0$ et $y\ge0$ : $200+7k\ge0\iff k\ge-28{,}5$ et
$-300-11k\ge0\iff k\le-27{,}27$. Le seul entier est $k=-28$, qui donne
$x=200-196=4$ et $y=-300+308=8$. Il faut donc \textbf{$4$ grands cartons et $8$ petits cartons}
($4\times11+8\times7=44+56=100$).

\noindent\textbf{Exercice 3.}
1) En $-\infty$ : $x\to-\infty$ et $\mathrm{e}^{-x}\to+\infty$, donc $f(x)\to-\infty$.
En $+\infty$ : $f(x)=\frac{x}{\mathrm{e}^{x}}\to0^+$ par croissances comparées : la droite
$y=0$ est asymptote horizontale en $+\infty$.
2) $f'(x)=1\cdot\mathrm{e}^{-x}+x\cdot(-\mathrm{e}^{-x})=(1-x)\mathrm{e}^{-x}$. Comme
$\mathrm{e}^{-x}>0$, $f'(x)>0\iff x<1$ : $f$ croît sur $]-\infty;1]$, décroît sur
$[1;+\infty[$, avec un maximum $f(1)=\mathrm{e}^{-1}=\frac1{\mathrm{e}}$.
3) IPP avec $u=x$, $v'=\mathrm{e}^{-x}$, donc $u'=1$, $v=-\mathrm{e}^{-x}$ :
$\int_0^{\lambda}x\,\mathrm{e}^{-x}\,\dd x=\big[-x\mathrm{e}^{-x}\big]_0^{\lambda}+\int_0^{\lambda}\mathrm{e}^{-x}\,\dd x
=-\lambda\mathrm{e}^{-\lambda}+\big[-\mathrm{e}^{-x}\big]_0^{\lambda}
=-\lambda\mathrm{e}^{-\lambda}+(1-\mathrm{e}^{-\lambda})=1-(\lambda+1)\mathrm{e}^{-\lambda}$.
4) Sur $[0;1]$, $f(x)=x\mathrm{e}^{-x}\ge0$. L'aire en unités d'aire vaut, avec $\lambda=1$,
$\int_0^1 f=1-2\mathrm{e}^{-1}$. L'unité étant \SI{2}{\centi\metre}, $1$ u.a. $=2\times2=\SI{4}{\centi\metre\squared}$,
donc $\mathcal{A}=4\,(1-2\mathrm{e}^{-1})\approx4\times0{,}264=\SI{1,06}{\centi\metre\squared}$.

\noindent\textbf{Exercice 4.}
1) La somme des coefficients est $2+1+1=4\ne0$, donc $G$ existe et
$\overrightarrow{AG}=\frac14\big(2\overrightarrow{AA}+\overrightarrow{AB}+\overrightarrow{AC}\big)
=\frac14(\overrightarrow{AB}+\overrightarrow{AC})$. En coordonnées,
$G=\frac14(2A+B+C)=\frac14\big((0,0)\!\cdot\!2+(4,0)+(0,4)\big)=\frac14(4,4)=(1\,;\,1)$.
2) Pour tout $M$ et tout point $P$, $\overrightarrow{PM}=\overrightarrow{PG}+\overrightarrow{GM}$. En
développant $MA^2=\overrightarrow{MA}\cdot\overrightarrow{MA}$ avec
$\overrightarrow{MA}=\overrightarrow{MG}+\overrightarrow{GA}$ (idem pour $B$, $C$) et en
utilisant $2\overrightarrow{GA}+\overrightarrow{GB}+\overrightarrow{GC}=\vec0$ (propriété du barycentre) :
$f(M)=2MA^2+MB^2+MC^2=(2{+}1{+}1)MG^2+\big(2GA^2+GB^2+GC^2\big)=4MG^2+(2GA^2+GB^2+GC^2)$.
Or $GA^2=1^2+1^2=2$, $GB^2=(4-1)^2+1^2=10$, $GC^2=1^2+(4-1)^2=10$, donc
$2GA^2+GB^2+GC^2=4+10+10=24$. D'où $f(M)=4\,MG^2+24$.
3) $f(M)=40\iff4MG^2+24=40\iff MG^2=4\iff MG=2$. L'ensemble $(\Gamma)$ est le
\textbf{cercle de centre $G(1\,;\,1)$ et de rayon $2$}.
4) $AG=\sqrt{1^2+1^2}=\sqrt2\ne2$, donc $A\notin(\Gamma)$ (on a $f(A)=4\times2+24=32\ne40$).

\smallskip
\noindent\textbf{Partie B — Situation.}
\emph{Tâche 1.} Le coût moyen est $C_m(x)=\dfrac{x^2+3600}{x}=x+\dfrac{3600}{x}$ pour $x>0$.
Alors $C_m'(x)=1-\dfrac{3600}{x^2}=\dfrac{x^2-3600}{x^2}$, qui s'annule en $x=60$ (positif),
avec $C_m'<0$ sur $]0;60[$ et $C_m'>0$ sur $]60;+\infty[$ : minimum en $x=60$. Il faut donc
produire $60$ centaines de baguettes (soit $6000$ baguettes) ; le coût moyen minimal est
$C_m(60)=60+\frac{3600}{60}=60+60=120$ milliers de francs CFA par centaine.

\emph{Tâche 2.} À intérêts composés, le capital au bout de $n$ années est
$C_n=500\,000\times1{,}06^{\,n}$. On cherche le plus petit $n$ tel que $C_n\ge2\times500\,000$,
c'est-à-dire $1{,}06^{\,n}\ge2$, soit $n\ge\dfrac{\ln2}{\ln1{,}06}\approx11{,}9$. Le capital
aura donc au moins doublé au bout de \textbf{$12$ ans} ($1{,}06^{12}\approx2{,}01$).

\emph{Tâche 3.} Chaque paquet est défectueux avec la probabilité $0{,}04$, indépendamment
des autres ; le nombre $Z$ de paquets défectueux parmi $5$ suit la loi binomiale
$\mathcal{B}(5\,;\,0{,}04)$. Donc
$P(Z\ge1)=1-P(Z{=}0)=1-(0{,}96)^5\approx1-0{,}815=\mathbf{0{,}185}$, soit environ \SI{18,5}{\percent}.
